\section{Script language}
\label{sec:language}

In this section we give a more detailed description of the script language.
We describe the types in Section~\ref{ssec:types}.  We describe expressions in
Section~\ref{ssec:expressions}.  We describe declarations of values and
functions in Section~\ref{ssec:declarations}.  We then describe cell write
commands in Section~\ref{ssec:cell-writes}, and |for| loop iterations in
Section~\ref{ssec:for-loops}. 

%%%%% 

\subsection{Types}
\label{ssec:types}

The script language contains the standard types |Int| of (32-bit) integers,
|Float| of (32-bit) floating point numbers, |Boolean| of booleans (with values
|true| and |false|), and |String| of strings.  

These four types (and no others) can be entered in cells of the spreadsheet.
The application decides the type of each value based on how it is entered,
mostly in the obvious way; for example, if a value looks like an integer, it
is given the type |Int|.  In most cases, a |String| can be entered by simply
typing the text, or can (optionally) be enclosed in quotation marks.  However,
in the case that a |String| is required that looks like a value of another
type, it \emph{must} be entered inside quotation marks (e.g.~|"3"|).  A
literal quotation mark can be entered as \SCALA{\\"} (e.g.~\SCALA{"He said
  \\"Hello\\""}).

The types |Row| and |Column| represent rows and columns, respectively. 

The type |Unit| is a type that contains a single value, called the unit value,
denoted~|()|.  This is the return type for functions that do not return a
proper result, but (typically) just perform some calculations and write into
cells. 

If |T| is a type, then |List[T]| represents the type of lists whose elements
are from~|T|.

If |T1| and |T2| are types, then |(T1,T2)| is the type of pairs $(x,y)$ with
$x \in \sm{T1}$ and $y \in \sm{T2}$.  This extends in the obvious way to
tuples of size up to four.

If |T1| and |T2| are types, then |T1 => T2| is the type of functions from |T1|
to~|T2|.  Note that the constructor ``|=>|'' associates to the right, so
\SCALA{T1 => T2 => T3} means \SCALA{T1 => (T2 => T3)}.



%%%%%%%%%%%%%%%%%%%%%%%%%%%%%%%%%%%%%%%%%%%%%%%%%%%%%%%

\subsection{Expressions}
\label{ssec:expressions}

\subsubsection{Standard expressions}

The script language contains the normal arithmetic operators, |+|, |-|, |*|
and~|/|, over |Int| and |Float| (over |Int|, |/| performs rounding), and
\SCALA{\%} over |Int|.  The language also contains the normal equality (|==|),
inequality (|!=|), and order (|<|, |<=|, |>|, |>=|) operators.   Note
that |Int| and |Float| are distinct types, so can't (currently) be combined
using these operators.  Unary negation can be written using |-|; note that the
expression might need to be included in parentheses when the |-| could be
interpreted as a binary operator, for example,
\begin{scala}
  { val x = 2
    (-1)
  }
\end{scala}
%
The built-in functions |toFloat: Int => Float| and
|toInt: Float => Int| can be used to convert between these types.

The language includes the normal boolean binary operators, |&&| and
\SCALA{\|\|}.  Boolean negation is written using the function |not|.

Conditional expressions are written using C/Java/Scala-style syntax: 
\begin{scala}
  if(£\syntax{test}£) £\syntax{expression}£ else £\syntax{expression}£
\end{scala}
%
The two branches must have the same type, and this is the overall type of the
conditional expression. 

Functional applications can be written using parentheses, in the normal way,
as in Figure~\ref{fig:example-script}.  However, as in Haskell, the
parentheses can be omitted for a function of one argument, where the actual
parameter does not logically need parentheses: for example, |f 3|, or \SCALA{g
  #A3}.  Function application associates to the left, and binds tighter than
infix operators; for example, \SCALA{f g true + h \#A} represents
\SCALA{((f(g))(true)) + (h(\#A))}.

String literals are written inside double quotation marks, as in most other
languages, e.g.~|"Hello"|.  Special characters can be written using
\SCALA{\\}, for example, \SCALA{"He said \\"Hello\\".\\n"}.

As noted in Section~\ref{ssec:types}, the unit value is denoted |()|. 

Tuple expressions are written inside parentheses, using commas, for example
\SCALA{(x+y, 3.6, x>y)}.  Tuples must contain at least two, and at most four
components.  When a function is applied to a tuple expression, an extra set of
parentheses is required, for example |f((x,y))| (since the expression |f(x,y)|
would imply that |f| takes two arguments, rather than a single tuple
argument).  The functions |get1|, |get2|, |get3|, and |get4| extract the
corresponding component from a tuple; for example $\sm{get2((2,4,6,8))} =
\sm{4}$.

A \emph{block expression} consists of a sequence of declarations, followed by
an expression, separated by semicolons or newlines, and enclosed in braces;
for example,
\begin{scala}
  { def square(x: Int) = x*x; val y2 = square y; val z2 = square z; y2+z2 }
\end{scala}
The declarations are evaluated in order; then the value of the final
expression is the value of the block expression.

%%%%%

\subsubsection{List expressions}

List literals are written in square brackets, for example |[1,2,3]| or |[]|.
All the elements of the list must be of the same type.  With the empty list,
it is sometimes necessary to give the type, to help the type checker, for
example |[]: List[Int]|.

List concatenation is written using |::|, pronounced ``cons''.  This operator
takes an element on the left and a list on the right, and returns their
concatenation; for example $\sm{3 :: [4,5]} = \sm{[3,4,5]}$.  The function
|head| returns the first element of a non-empty list; for example $\sm{head
  [3,4,5]} = \sm{3}$.  The function |tail| returns the list containing all
except the first element of a non-empty list; for example $\sm{tail [3,4,5]} =
\sm{[4,5]}$.  The function |isEmpty| tests whether a list is empty. 

Lists of consecutive |Int|s can be constructed using the infix operators |to|
and |until|.  The list |m to n| produces the list from~|m| to~|n| inclusive;
for example, $\sm{3 to 5} = \sm{[3,4,5]}$.  The list |m until n| is similar,
but does not include~|n|; for example, $\sm{3 until 5} = \sm{[3,4]}$.  These
operators can also be applied to row and column values; see below. 

List comprehensions provide a way to define a list, in a similar style to
mathematical set comprehensions.  The expression \SCALA{[e \| x <- xs]}
generates the list containing the value of~|e| for each~|x| in~|xs|.  For
example,
%
\begin{eqnarray*}
\mbox{\SCALA{[x*x \| x <- [3,4,5]]}} & = & \sm{[9,16,25]}.
\end{eqnarray*}
%
A term of the form |x <- xs| is called a \emph{generator}.  A list
comprehension may have multiple generators, separated by commas; an earlier
generator is iterated over more slowly than a later generator.  For example,
\begin{eqnarray*}
\mbox{\SCALA{[x*y | x <- [2,3], y <- [4,5]]}} & = & \sm{[8,10,12,15]}.
\end{eqnarray*}
%
A list comprehension may also have one or more boolean filters; a term is
produced only for values of the variables that satisfy the filter.  For
example, 
%
\begin{eqnarray*}
\mbox{\SCALA{[x*y | x <- [2,3], y <- [4,5], x+y != 7]}} & = & \sm{[8,15]}.
\end{eqnarray*}

A generator may perform pattern matching over tuples.  A generator of the form
|(x,y) <- pairs| expects |pairs| to be a list of pairs, and binds the
names~|x| and~|y| to the two components of each pair in turn.  For example,
%
\begin{eqnarray*}
\mbox{\SCALA{[x*y | (x,y) <- [(2,4), (3,5)]]}} & = & \sm{[8,15]}.
\end{eqnarray*}
%
This extends to tuples with more than two elements, and to nested tuples, in
the obvious way.

%%%%%

\subsubsection{Row, column, and cell expressions}

Row and column literals are written using |#|, e.g.~|#A|, |#B|, etc., and
|#0|, |#1|, etc.  A row or column value can have an |Int| added to it or
subtracted from it, with the obvious effect.  For example, $\sm{\#C} + 2 =
\sm{\#E}$, and $\sm{\#C} - 2 = \sm{\#A}$.  Arithmetic underflow, e.g.~from
$\sm{\#C} - 3$, gives an error.  Similarly, a row value can have another row
value subtracted from it; and a column value can have another column value
subtracted from it.  For example, $\sm{\#5} - \sm{\#2} = \sm{\#E} - \sm{\#B} =
3$. 

Lists of row and column values can be formed using the |to| and |until|
operators.  For example \SCALA{#B to #D} $=$ \SCALA{[#B,#C,#D]}, and \SCALA{#2
  until #4} $=$ \SCALA{[#2,#3]}.


If |c| and |r| are a |Row| and |Column|, respectively, then |Cell(c,r)| is a
\emph{cell expression}: it represents the value in the corresponding cell of
the spreadsheet.  This can be abbreviated in the case of row and column
literals; for example \SCALA{\#D3} is shorthand for \SCALA{Cell(\#D, \#3)}.

The type checker needs to be able to deduce a type, or a set of possible
types, for each cell expression.  When the script is executed, the application
checks if the corresponding cell in the spreadsheet has an expected type, and
signals an error otherwise.

One option is to provide an explicit type for a cell expression.  For example,
\SCALA{\#D3: Int} gives the type |Int| to the cell expression.

Alternatively, the type can often be deduced automatically if a function is
applied to the cell expression: this implies that its type must match the
parameter type for the function.  For example, in line~\ref{line:write-facts}
of Figure~\ref{fig:example-script}, the expression |fact(Cell(In,r))| implies
that the cell expression must have type |Int|, since |fact| takes a parameter
of type~|Int|.

Similarly, the type deduction can sometimes be made when a binary operator is
applied to a cell expression.  For example, consider the expression \SCALA{1 +
  #A3}; the type checker works from left to right, so it can deduce that the
cell expression must have type |Int|.  However, the explicit type is necessary
in the similar expression \SCALA{#A3:Int + 1}, because |+| is overloaded, and
the type checker works from left to right.

Finally, a |match| expression allows the cell to have one of several different
types, and for the value of the expression to depend on that type.  An example
was given on line~\ref{line:cell-match} of Figure~\ref{fig:example-script}.
In general, a |match| expression takes the form
\begin{scala}
  £\syntax{cell expression}£ match{ £\syntax{cases}£ } 
\end{scala}
where \syntax{cases} is a sequence of cases, separated by semicolons or
newlines.  Each case is of the form
\begin{scala}
  case £\syntax{pattern}£ => £\syntax{expression}£
\end{scala}
%
Each pattern takes one of the following forms.
\begin{itemize}
\item |Empty|, which matches an empty cell;

\item \syntax{name}|: |\syntax{type}, where \syntax{type} is one of |Int|,
  |Float|, |String|, or |Boolean|; this matches a value of the given type, and
  binds the name to that value in the subsequent expression;
 
\item |_|\,|: |\syntax{type}, which is similar to the previous form, except no
  name is bound;

\item |_|\,, which matches any type.
\end{itemize}
%
In a |match| expression, the cases are tried in the order given; the first
case that matches is selected, and the corresponding expression is evaluated.
If no case matches, an error is signalled.  Within a |match| expression, all
the expressions on the right-hand side of cases must have the same type.  

